\documentclass[11pt]{article}

\usepackage[margin=1in]{geometry}
\usepackage[T1]{fontenc}
\usepackage{lmodern}
\usepackage{enumitem}
\usepackage[hidelinks]{hyperref}

\setlist[itemize]{leftmargin=*, itemsep=4pt, topsep=4pt}
\setlength{\parindent}{0pt}
\setlength{\parskip}{6pt}

\begin{document}

\section*{Timeline}

\subsection*{September (Mathematical Formalization)}
\begin{itemize}
  \item Get a basic handle on all of the concepts of both questions
  \item Check for novelty and feasibility
  \item Research deeper into the two potential topics and decide on one before the end of September
\end{itemize}

\subsection*{October (Base Infrastructure)}
\begin{itemize}
  \item Continue working on mathematical formalization as many parts of it will most likely still be incomplete by the end of September
  \item Have basic optimal transport (OT) framework running
  \item Add a variety of tests to the framework
  \item Begin to plan how to extend the framework to work with MOTs
\end{itemize}

\subsection*{November (Neural Solver Integration)}
\begin{itemize}
  \item Finish the extension of the framework so that it can handle MOTs
  \item Complete the implementation of the neural network (CINN/PINN)
  \item Development of the loss function
  \item Run extensive testing to validate correctness of code
\end{itemize}

\subsection*{December (Parameterization \& Initial Training)}
\begin{itemize}
  \item Look into methods to formally prove the correctness of code
  \item Start training / parameterizing of all neural networks, utilizing our continuity equations and martingale conditions
  \item If possible start to use standard 1D and 2D Black--Scholes dynamics as benchmark reference solutions before scaling to multi-asset continuum limits to verify implementation and direction
\end{itemize}

\end{document}
